%% ============================================================
%% 分片 25：原文第 9 节「't Hooft And Hecke Operators」（原文 glshort.tex 起于第 6755 行）
%% 由 tools/harvmac2tex.py 从 arXiv 源码（hep-th/0604151v3）自动生成。
%% 只译英文散文；公式、编号、\label、\cite 一律不改。
%% ============================================================

\subsection{The Affine Grassmannian}\label{sec:affgrass}

The space of Hecke modifications associated with an 't Hooft
operator can be described as a finite-dimensional subvariety in a
certain infinite-dimensional variety called the affine
Grassmannian. This description enables one to construct a natural
compactification of the space of Hecke modifications.  We will not
really use this framework in the present paper, but we explain it
here because it gives a powerful way of understanding some things
and is the framework in which the subject is usually placed
mathematically. The affine Grassmannian is familiar to physicists
primarily for its applications to two-dimensional conformal field
theory \cite{ref122}.

For the sake of clarity, let us consider the case $G=U(N)$. In
this case, an 't Hooft operator is specified by a set of integers
$m_1,\ldots,m_N$ modulo permutations. Since the problem is local,
we can fix the curve $C$ to be $\Bbb{CP}^1$ and take $E_-$ to be
trivial.  A Hecke modification $E$ of $E_-$ at a point $p_0\in C$
is a bundle that is endowed with a chosen isomorphism
$\sigma:E\cong E_-$ outside of $p_0$; $\sigma$ may not extend over
$p_0$.

It is natural to regard the space of Hecke modifications for fixed
$m_1,\ldots,m_N$ as a subspace of the space of pairs $(E,\sigma)$,
where $E$ is a holomorphic vector bundle on $C$ and $\sigma$ is
its trivialization outside $p_0$. Let us recall an explicit
description of the space of such pairs $(E,\sigma)$, following
chapter 8 of \cite{ref122}.

It is convenient to think of $C=\Bbb{CP}^1$ as a one-point
compactification of the complex $z$-plane $\CC$ and to set $p_0$
to be the point $z=0$. We also let % 式 (9.28)
\setcounter{equation}{27}
\begin{equation}\label{eq:jmj}
U_\infty=\Bbb{CP}^1-\{0\},\quad U_0=\Bbb{CP}^1-\{\infty\} .
\end{equation}
 We
are given a trivialization $\sigma$ of $E$ over $U_\infty$.
Although $\sigma$ may not extend over $U_0$, we can pick a
trivialization $\sigma'$ of $E$ over $U_0$. Over $U_0\bigcap
U_\infty\simeq\CC^*$, $\sigma $ and $\sigma'$ are related by a
$GL(N,\CC)$ gauge transformation $g(z)$. This is a
$GL(N,\CC)$-valued function whose entries are holomorphic on
$\CC^*$ but may have poles at $0$ and $\infty$. ($\sigma'$ can
always be chosen so that the singularities of $g(z)$ are poles.)
Let us denote by ${\cal X}$ the ring of functions holomorphic on
$\CC^*$ and having a meromorphic extension to $\Bbb{CP}^1$; then
the group of holomorphic gauge transformations on $\CC^*$ can be
identified with $GL(N,{\cal X})$, and $g(z)$ is an element of this
group.

If we change $\sigma'$ by a gauge transformation $h_0(z)$ which is
holomorphic throughout $U_0$, then $g(z)$ is replaced by
$g(z)h_0^{-1}(z)$.  $h_0(z)$ takes values in $GL(N,\cO)$, the
group of invertible matrices whose entries are holomorphic
functions on $U_0\simeq \CC$. The set of holomorphic vector
bundles of rank $N$ over $\Bbb{CP}^1$ equipped with a fixed
holomorphic trivialization $\sigma$ outside $p_0$ is isomorphic to
the quotient $GL(N,{\cal X})/GL(N,\cO)$. The latter space is known
as the affine Grassmannian ${\rm Gr}_N$ for the group $GL(N)$.
Another name for it is the loop Grassmannian, because it is
isomorphic \cite{ref122}\ to the space of based loops in $U(N)$.

The definition of the affine Grassmannian admits slight
variations. For example, since any function in ${\cal X}$ can be
expanded in a Laurent series, one can embed ${\cal X}$ into the
ring $\CC((z))$ of formal Laurent series. Similarly, one can embed
the ring of holomorphic functions $\cO$ into the ring $\CC[[z]]$
of formal power series. One can show that the quotient
$GL(N,\CC((z)))/GL(N,\CC[[z]])$ is isomorphic to the affine
Grassmannian ${\rm Gr}_N$. Intuitively, this happens because given
any two elements of $GL(N,{\cal X})$, one can tell whether they
are in the same $GL(N,\cO)$ orbit by studying a finite number of
terms in their Laurent expansions, and therefore it is immaterial
whether the Laurent series have a nonzero region of convergence.

Now we have to identify those points in ${\rm Gr}_N$ which can be
obtained from the trivial vector bundle by a Hecke modification of
weight $^L\negthinspace w=(m_1,\ldots,m_N)$. A Hecke modification
of the trivial vector bundle with this weight is obtained by
choosing a trivialization of the trivial vector bundle over $U_0$
by $N$ linearly independent holomorphic sections $f_1,f_2,\dots,
f_N$, and declaring that $E_+$ is generated over $U_0$ by sections
$s_1,\ldots,s_N$ which upon restriction to $U_\infty\bigcap U_0$
are given by % 式 (9.29)
\setcounter{equation}{28}
\begin{equation}\label{eq:mcop}
z^{-m_j} f_j,\quad j=1,\ldots,N.
\end{equation}
 Here we
made use of the fact that on $U_\infty$ we are given an
isomorphism between $E_+$ and the trivial vector bundle. Thus if
$f_j=(f_j^1,f_j^2,\dots,f_j^N)$, then the matrix $g(z)$
corresponding to $E_+$ is given by % 式 (9.30)
\setcounter{equation}{29}
\begin{equation}\label{eq:xcoxo}
g^i_j(z)=z^{-m_j}
f^i_j(z).
\end{equation}
 The simplest choice is $f^i_j(z)=\delta^i_j$; all
other choices can be obtained from this one by acting on $g(z)$
from the left by an element of $GL(N,\cO)$. We conclude that the
space of Hecke modifications ${{\EUY}}({}^L\negthinspace w;p_0)$
is the orbit of the point % 式 (9.31)
\setcounter{equation}{30}
\begin{equation}\label{eq:gspec}
g^i_j(z)=z^{-m_j}
\delta^i_j
\end{equation}
 in the affine Grassmannian under the left action of
the group $GL(N,\cO)$. One can replace ${\cal X}$ with $\CC((z))$
and $\cO$ with $\CC[[z]]$ throughout and get an equivalent result.
Note that $g(z)$ given by \eqref{eq:gspec}\ is a homomorphism of $\CC^*$ to
$GL(N,\CC)$.  In fact, it is the complexification of the
homomorphism % 式 (9.32)
\setcounter{equation}{31}
\begin{equation}\label{eq:xxconon}
\rho: U(1)\to U(N),\quad \rho:
e^{i\alpha}\to \left(\begin{matrix}e^{i m_1\alpha} & & \ldots & \\ &
e^{i m_2\alpha} & \ldots &  \\  & \vdots & \ldots & \\ & & &
e^{i m_N\alpha}\\\end{matrix}\right)
\end{equation}
 which enters into the definition of
the 't Hooft operator $T(m_1,\ldots,m_N)$.

We have associated to each set of integers $m_1,\ldots,m_N$ an
orbit of $GL(N,\cO)$ in ${\rm Gr}_N$. It turns out that all
$GL(N,\cO)$ orbits in ${\rm Gr}_N$ are obtained in this way \cite{ref122} .
Thus ${\rm Gr}_N$ is stratified by spaces
${{\EUY}}(m_1,\ldots,m_N)$, which are called Schubert cells.
Equivalently, by applying Hecke modifications to the trivial
vector bundle, one can obtain an arbitrary holomorphic vector
bundle on $\Bbb{CP}^1$ with an arbitrary trivialization on
$U_\infty$.

The examples of spaces of Hecke modifications discussed in section
第~\ref{sec:spaceheck}~节\ suggest that these spaces are always
finite-dimensional. This can be shown in general as follows. Let
${\rm Gr}_N(i)$ be the subset of ${\rm Gr}_N$ defined by the
condition that $g(z)$ has at $z=0$ a pole of order not higher than
$i$. The variety ${\rm Gr}_N(i)$ is finite-dimensional,
$GL(N,\cO)$-invariant, and compact, and any point in ${\rm Gr}_N$
belongs to ${\rm Gr}_N(i)$ for some $i$. This implies that any
orbit of $GL(N,\cO)$ belongs to some ${\rm Gr}_N(i)$ and therefore
is finite-dimensional.

It may be helpful to mention at this point that the
infinite-dimensional space ${\rm Gr}_N$ has another stratification
with strata labelled by sets of integers $k_1,\ldots,k_N$. To
define this stratification, we recall that by a theorem of
Grothendieck any holomorphic vector bundle of rank $N$ on
$\Bbb{CP}^1$ is isomorphic to % 式 (9.33)
\setcounter{equation}{32}
\begin{equation}\label{eq:Groth}
\oplus_{i=1}^N
\cO(p_0)^{k_i}
\end{equation}
 for some integers $k_1,\ldots,k_N$. We define the
stratum in ${\rm Gr}_N$ corresponding to $k_1,\ldots,k_N$ as the
set of those pairs $(E,\sigma)$ where $E$ is isomorphic to \eqref{eq:Groth}.
In other words, a stratum is obtained by fixing $E$ and varying
the trivialization $\sigma$. Obviously, the strata are
infinite-dimensional in this case, in contrast with the
stratification given by ${{\EUY}}(m_1,\ldots,m_N)$. The relation
between the two stratifications of ${\rm Gr}_N$ is studied in \cite{ref122}.

As we have seen, spaces of Hecke modifications or Schubert cells
are noncompact, in general. A natural way to compactify them is to
consider the closure of the corresponding orbits in ${\rm
Gr}_N(i)$ for sufficiently large $i$. For a given $\rho: U(1)\to
U(N)$, the closure of the orbit $C_\rho$ is in general a singular
variety which  is a union of $C_\rho$ and a finite number of
orbits of lower dimension. It is called a Schubert cycle.  As
discussed in section 第~\ref{sec:spaceheck}~节, the structure of the closure of
$C_\rho$ reflects the ``mixing'' between 't Hooft operators with
different $\rho$ but the same topological type.  We also will
describe this process in terms of monopole bubbling in section
第~\ref{sec:monbub}~节.

The construction of ${\rm Gr}$ and the description of the spaces
of Hecke modifications as subvarieties in ${\rm Gr}$ can be
generalized to other gauge groups. For any simple compact Lie
group $G$, with complexification $G_{\Bbb{C}}$, one defines ${\rm
Gr}_G$ as the quotient % 式 (9.34)
\setcounter{equation}{33}
\begin{equation}\label{eq:hxonoh}
{\rm Gr}_G=G_{\Bbb{C}}({\cal
X})/G_{\Bbb{C}}(\cO).
\end{equation}
 An 't Hooft operator is parametrized by
the conjugacy class of a homomorphism $\rho: U(1)\to G$, which can
be analytically continued to a homomorphism $\rho_\CC: \CC^*\to
G_{\Bbb{C}}$. Obviously, $\rho_\CC$ defines a point on ${\rm
Gr}_G$, and the orbit of this point under the left action of
$G_{\Bbb{C}}(\cO)$ depends only on the conjugacy class of $\rho$.
This orbit is the Schubert cell or space of Hecke modifications
${{\EUY}}(\rho)$.
